
\documentclass{article}
\usepackage{colortbl}
\usepackage{makecell}
\usepackage{multirow}
\usepackage{supertabular}

\begin{document}

\newcounter{utterance}

\centering \large Interaction Transcript for game `referencegame', experiment `number\_grids', episode 17 with reflexion{-}Qwen3.5{-}4B.
\vspace{24pt}

{ \footnotesize  \setcounter{utterance}{1}
\setlength{\tabcolsep}{0pt}
\begin{supertabular}{c@{$\;$}|p{.15\linewidth}@{}p{.15\linewidth}p{.15\linewidth}p{.15\linewidth}p{.15\linewidth}p{.15\linewidth}}
    \# & $\;$A & \multicolumn{4}{c}{Game Master} & $\;\:$B\\
    \hline

    \theutterance \stepcounter{utterance}  
    & & \multicolumn{4}{p{0.6\linewidth}}{
        \cellcolor[rgb]{0.9,0.9,0.9}{
            \makecell[{{p{\linewidth}}}]{
                \texttt{\tiny{[P1$\langle$GM]}}
                \texttt{You are given three grids, where each of them is 5 by 5 in size.} \\
\texttt{Grids have empty cells marked with "▢" and filled cells marked with "X".} \\
\texttt{Your task is to generate a referring expression that best describes the target grid while distinguishing it from the two other distractor grids.} \\
\texttt{The first grid is the target grid, and the following two grids are the distractors.} \\
\\ 
\texttt{Target grid:} \\
\\ 
\texttt{X X X X □} \\
\texttt{X □ □ □ □} \\
\texttt{X X X X □} \\
\texttt{X □ □ □ X} \\
\texttt{X X X X □} \\
\\ 
\texttt{Distractor grid 1:} \\
\\ 
\texttt{X X X X □} \\
\texttt{X □ □ □ □} \\
\texttt{X X X X □} \\
\texttt{□ □ □ X □} \\
\texttt{X X X X □} \\
\\ 
\texttt{Distractor grid 2:} \\
\\ 
\texttt{□ X X X X} \\
\texttt{□ X □ □ X} \\
\texttt{□ X X X X} \\
\texttt{□ X □ □ X} \\
\texttt{□ X X X X} \\
\\ 
\texttt{Instruction: Describe the target grid.} \\
\texttt{Generate the referring expression starting with the tag "Expression: " for the given target grid. Omit any other text.} \\
            }
        }
    }
    & \\ \\

    \theutterance \stepcounter{utterance}  
    & \multicolumn{4}{p{0.6\linewidth}}{
        \cellcolor[rgb]{0.9,0.9,0.9}{
            \makecell[{{p{\linewidth}}}]{
                \texttt{\tiny{[P1$\rangle$GM]}}
                \texttt{Expression: The grid with a solid vertical bar of five X's in the first column, a solid horizontal bar of four X's in the first row, and a solid horizontal bar of four X's in the third row, with the bottom{-}right corner cell filled with an X.} \\
            }
        }
    }
    & & \\ \\

    \theutterance \stepcounter{utterance}  
    & & & \multicolumn{2}{p{0.3\linewidth}}{
        \cellcolor[rgb]{0.9,0.9,0.9}{
            \makecell[{{p{\linewidth}}}]{
                \texttt{\tiny{[GM$|$GM]}}
                \texttt{Expression: The grid with a solid vertical bar of five X's in the first column, a solid horizontal bar of four X's in the first row, and a solid horizontal bar of four X's in the third row, with the bottom{-}right corner cell filled with an X.} \\
            }
        }
    }
    & & \\ \\

    \theutterance \stepcounter{utterance}  
    & & \multicolumn{4}{p{0.6\linewidth}}{
        \cellcolor[rgb]{0.9,0.9,0.9}{
            \makecell[{{p{\linewidth}}}]{
                \texttt{\tiny{[GM$\rangle$P2]}}
                \texttt{You are given three grids, where each of them is 5 by 5 in size.} \\
\texttt{Grids have empty cells marked with "▢" and filled cells marked with "X".} \\
\texttt{You are also given a referring expression that describes one of the given grids.} \\
\texttt{Your task is to select the grid that matches the given referring expression.} \\
\texttt{Generate only the number (in text) of the grid that matches the given expression by selecting first, second, or third.} \\
\\ 
\texttt{First:} \\
\\ 
\texttt{□ X X X X} \\
\texttt{□ X □ □ X} \\
\texttt{□ X X X X} \\
\texttt{□ X □ □ X} \\
\texttt{□ X X X X} \\
\\ 
\texttt{Second:} \\
\\ 
\texttt{X X X X □} \\
\texttt{X □ □ □ □} \\
\texttt{X X X X □} \\
\texttt{□ □ □ X □} \\
\texttt{X X X X □} \\
\\ 
\texttt{Third:} \\
\\ 
\texttt{X X X X □} \\
\texttt{X □ □ □ □} \\
\texttt{X X X X □} \\
\texttt{X □ □ □ X} \\
\texttt{X X X X □} \\
\\ 
\texttt{Expression: The grid with a solid vertical bar of five X's in the first column, a solid horizontal bar of four X's in the first row, and a solid horizontal bar of four X's in the third row, with the bottom{-}right corner cell filled with an X.} \\
\texttt{Question: Which grid does the expression refer to?} \\
\texttt{Start with the tag "Answer: ", followed by your selection. Omit any other text.} \\
            }
        }
    }
    & \\ \\

    \theutterance \stepcounter{utterance}  
    & & & \multicolumn{4}{p{0.6\linewidth}}{
        \cellcolor[rgb]{0.9,0.9,0.9}{
            \makecell[{{p{\linewidth}}}]{
                \texttt{\tiny{[GM$\langle$P2]}}
                \texttt{Answer: Second} \\
            }
        }
    }
     \\ \\

    \theutterance \stepcounter{utterance}  
    & & & \multicolumn{2}{p{0.3\linewidth}}{
        \cellcolor[rgb]{0.9,0.9,0.9}{
            \makecell[{{p{\linewidth}}}]{
                \texttt{\tiny{[GM$|$GM]}}
                \texttt{second} \\
            }
        }
    }
    & & \\ \\

    \theutterance \stepcounter{utterance}  
    & \multicolumn{4}{p{0.6\linewidth}}{
        \cellcolor[rgb]{0.95,0.95,0.95}{
            \makecell[{{p{\linewidth}}}]{
                \texttt{\tiny{[P2$\rightarrow$P1]}}
                \texttt{Post{-}episode Reflexion (not sent to GM):} \\
\texttt{New plan: I failed because my description was too generic and failed to uniquely distinguish the target grid from the distractors. Specifically, I described the common features (solid bars in the first column and first/third rows) which are shared by the target and distractor 1, and I missed the critical difference: the target grid has an 'X' in the bottom{-}right corner (row 5, column 5), while distractor 1 has a '▢' there. Next time, I will construct the expression by first stating the unique feature that separates the target from the others (the bottom{-}right 'X'), then describing the shared structural elements (the solid bars) to confirm the overall shape, ensuring the description is logically sufficient to identify the target among the three options.} \\
            }
        }
    }
    & & \\ \\

    \theutterance \stepcounter{utterance}  
    & & & \multicolumn{4}{p{0.6\linewidth}}{
        \cellcolor[rgb]{0.95,0.95,0.95}{
            \makecell[{{p{\linewidth}}}]{
                \texttt{\tiny{[P1$\rightarrow$P2]}}
                \texttt{Post{-}episode Reflexion (not sent to GM):} \\
\texttt{New plan: I incorrectly identified the target grid by relying on a partial match of the first column and first row, failing to verify the specific constraints for the third row and the bottom{-}right corner. The expression requires a solid horizontal bar of four X's in the third row (meaning the bottom{-}right cell must be empty, as it is the fifth cell), but the expression also explicitly states the bottom{-}right corner must be filled with an X. This is a contradiction in the description or my reading of the grid dimensions. Re{-}evaluating the grids: Grid 1 has no X in the first row's last column. Grid 2 has an X in the first row's last column, but its third row has an X in the fourth column and an empty fifth column (bottom{-}right is empty). Grid 3 has an X in the first row's last column, an X in the third row's fourth column, and an X in the bottom{-}right corner. Wait, the expression says "solid horizontal bar of four X's in the third row". In a 5x5 grid, a bar of four X's in the third row would occupy columns 1{-}4 or 2{-}5. If it's columns 1{-}4, the bottom{-}right (col 5) is empty. If it's columns 2{-}5, the first column of that row is empty. The expression also says "solid vertical bar of five X's in the first column". This implies column 1 is all X. Therefore, the third row must have X in column 1. If the bar is 4 X's starting at col 1, it covers cols 1{-}4. This leaves col 5 empty. But the expression says "bottom{-}right corner cell filled with an X". This is impossible if the third row has a solid bar of 4 X's starting from the first column (which is required by the vertical bar). Unless the "solid horizontal bar" is in columns 2{-}5? But then the first column of the third row would be empty, violating the vertical bar. I misread the grid structures or the expression's logical consistency. Let's re{-}read carefully: "solid vertical bar of five X's in the first column" {-}> Col 1 is all X. "solid horizontal bar of four X's in the first row" {-}> Row 1 has 4 X's. "solid horizontal bar of four X's in the third row" {-}> Row 3 has 4 X's. "bottom{-}right corner cell filled with an X" {-}> (5,5) is X.} \\
\texttt{Grid 1: Row 1 is empty except col 2. Fails.} \\
\texttt{Grid 2: Row 1 is X X X X □. (4 X's). Col 1 is X X X X X. (5 X's). Row 3 is X X X X □. (4 X's). Bottom{-}right} \\
            }
        }
    }
     \\ \\

\end{supertabular}
}

\end{document}
